% Blueprint content for YoungDiagram project
% Reference: Đoković's classification of nilpotent orbits in classical real Lie groups

\chapter{Genes and Chromosomes}

\section{Gene types and genes}

\begin{definition}[Gene type]
  \label{def:GeneType}
  \lean{GeneType}
  \leanok
  A \emph{gene type} is one of three polarization types:
  $\mathtt{NonPolarized}$, $\mathtt{Positive}$, or $\mathtt{Negative}$.
  There is a natural involution $\varepsilon \mapsto -\varepsilon$ that swaps
  $\mathtt{Positive}$ and $\mathtt{Negative}$ and fixes $\mathtt{NonPolarized}$.
\end{definition}

\begin{definition}[Gene]
  \label{def:Gene}
  \lean{Gene}
  \leanok
  \uses{def:GeneType}
  A \emph{gene} $g = (n, \varepsilon)$ consists of a \emph{rank} $n \in \mathbb{N}$
  and a \emph{type} $\varepsilon \in \mathtt{GeneType}$.
\end{definition}

\begin{definition}[Genes of a fixed rank]
  \label{def:Gene.ofRank}
  \lean{Gene.ofRank, Gene.ofRankAlt}
  \leanok
  \uses{def:Gene}
  For $n \ge 1$ and a type $\varepsilon$, $\mathtt{Gene.ofRank}\,n\,\varepsilon$
  is the one-gene chromosome consisting of the gene $(n,\varepsilon)$ with
  multiplicity one; for $n=0$ it is the zero chromosome.  The variant
  $\mathtt{Gene.ofRankAlt}$ uses the alternating type
  $(-1)^{n-1}\varepsilon$.
\end{definition}

\begin{definition}[Gene signature]
  \label{def:Gene.signature}
  \lean{Gene.signature}
  \leanok
  \uses{def:Gene}
  The \emph{signature} of a gene $g = (n, \varepsilon)$ is a pair
  $(a, b) \in \mathbb{Q} \times \mathbb{Q}$ defined by:
  \begin{itemize}
    \item If $\varepsilon = \mathtt{NonPolarized}$: $(a,b) = (n/2,\, n/2)$.
    \item If $\varepsilon = \mathtt{Positive}$: $(a,b) = (\lceil n/2 \rceil,\, \lfloor n/2 \rfloor)$.
    \item If $\varepsilon = \mathtt{Negative}$: $(a,b) = (\lfloor n/2 \rfloor,\, \lceil n/2 \rceil)$.
  \end{itemize}
\end{definition}

\begin{lemma}[Signature sum equals rank]
  \label{lem:Gene.signature_sum_eq_rank}
  \lean{Gene.signature_sum_eq_rank}
  \leanok
  \uses{def:Gene.signature}
  For any gene $g$, we have $a(g) + b(g) = \mathrm{rank}(g)$.
\end{lemma}

\section{Chromosomes}

\begin{definition}[Chromosome]
  \label{def:Chromosome}
  \lean{Chromosome}
  \leanok
  \uses{def:Gene}
  A \emph{chromosome} is a finitely supported function $X \colon \mathtt{Gene} \to \mathbb{N}$,
  i.e.\ a formal sum of genes with non-negative integer multiplicities.
  We write $\mathtt{Chromosome} = \mathtt{Gene} \to_0 \mathbb{N}$.
\end{definition}

\begin{definition}[Chromosome signature]
  \label{def:Chromosome.signature}
  \lean{Chromosome.signature}
  \leanok
  \uses{def:Chromosome, def:Gene.signature}
  The \emph{signature} of a chromosome $X$ is the additive extension:
  $\mathrm{sig}(X) = \sum_{g} X(g) \cdot \mathrm{sig}(g) \in \mathbb{Q} \times \mathbb{Q}$.
\end{definition}

\begin{lemma}[Signature of a basic gene]
  \label{lem:Chromosome.signature_ofRank}
  \lean{Chromosome.signature_ofRank, Chromosome.signature_ofRank_eq, Chromosome.signature_ofRank_eq'}
  \leanok
  \uses{def:Gene.ofRank, def:Chromosome.signature}
  The signature of $\mathtt{Gene.ofRank}\,n\,\varepsilon$ is the signature of
  the corresponding gene when $n\ne 0$, and zero when $n=0$.  For polarized
  types the file also records the one-step recurrence obtained by lowering the
  rank and changing, or preserving, the sign according to the parity.
\end{lemma}

\begin{definition}[Rank]
  \label{def:Chromosome.rank}
  \lean{Chromosome.rank}
  \leanok
  \uses{def:Chromosome}
  The \emph{rank} of a chromosome $X$ is $\mathrm{rank}(X) = \sum_g X(g) \cdot \mathrm{rank}(g)$.
\end{definition}

\begin{lemma}[Signature sum equals rank for chromosomes]
  \label{lem:Chromosome.signature_sum_eq_rank}
  \lean{Chromosome.signature_sum_eq_rank}
  \leanok
  \uses{def:Chromosome.signature, def:Chromosome.rank}
  For any chromosome $X$, the sum of the two signature components equals the rank:
  $\mathrm{sig}(X).1 + \mathrm{sig}(X).2 = \mathrm{rank}(X)$.
\end{lemma}

\begin{lemma}[Zero signature forces zero chromosome]
  \label{lem:Chromosome.signature_eq_zero}
  \lean{Chromosome.signature_eq_zero}
  \leanok
  \uses{def:Chromosome.signature, def:Chromosome.rank, lem:Chromosome.signature_sum_eq_rank}
  If $\mathrm{sig}(X)=0$, then $X=0$.
\end{lemma}

\begin{definition}[Dominance order]
  \label{def:Chromosome.Dominates}
  \lean{Chromosome.Dominates}
  \leanok
  \uses{def:Chromosome.signature, def:Chromosome.prime}
  For chromosomes $X$ and $Y$, we say $X$ \emph{dominates} $Y$ (written $Y \le X$) if
  $\mathrm{sig}(Y^{(k)}) \le \mathrm{sig}(X^{(k)})$ componentwise for all $k \ge 0$,
  where $X^{(k)}$ denotes the $k$-th derivative.
\end{definition}

\begin{lemma}[Dominance is compatible with addition]
  \label{lem:Chromosome.ordered_cancel}
  \lean{Chromosome.le_iff_dominates, Chromosome.sub_single_lt_sub_single}
  \leanok
  \uses{def:Chromosome.Dominates}
  The dominance preorder is an ordered cancel additive monoid structure on
  chromosomes.  In particular, subtracting the same present gene from two
  strictly ordered chromosomes preserves strict order.
\end{lemma}

\section{Prime and lift operations}

\begin{definition}[Prime (derivative)]
  \label{def:Chromosome.prime}
  \lean{Chromosome.prime}
  \leanok
  \uses{def:Chromosome}
  The \emph{prime} (derivative) operation $X \mapsto X'$ is the additive homomorphism
  that maps a gene $(n, \varepsilon)$ to the gene $(n-1,\varepsilon)$ when
  $n>1$, and sends rank-one genes to $0$.
\end{definition}

\begin{lemma}[Prime on generators and coefficients]
  \label{lem:Chromosome.prime_generators}
  \lean{Chromosome.prime_ofRank, Chromosome.prime_iterate_ofRank, Chromosome.prime_coeff}
  \leanok
  \uses{def:Gene.ofRank, def:Chromosome.prime}
  Prime lowers the rank of an $\mathtt{ofRank}$ generator by one, iterated
  prime lowers it by the number of iterations, and the coefficient of a gene
  $g$ in $X'$ is the coefficient of the rank-shifted gene $(g.rank+1,g.type)$
  in $X$.
\end{lemma}

\begin{definition}[Lift (antiderivative)]
  \label{def:Chromosome.lift}
  \lean{Chromosome.lift}
  \leanok
  \uses{def:Chromosome}
  The \emph{lift} operation is an additive homomorphism mapping a gene
  $(n,\varepsilon)$ to $(n+1,\varepsilon)$.
\end{definition}

\begin{lemma}[Prime-lift left inverse]
  \label{lem:prime_lift_leftInverse}
  \lean{Chromosome.prime_lift_leftInverse}
  \leanok
  \uses{def:Chromosome.prime, def:Chromosome.lift}
  The prime operation is a left inverse to lift: $(\mathrm{lift}\, X)' = X$ for all $X$.
\end{lemma}

\begin{lemma}[Max rank decreases under prime]
  \label{lem:maxRank_prime_lt}
  \lean{Chromosome.maxRank_prime_lt}
  \leanok
  \uses{def:Chromosome.prime}
  If $X \ne 0$, then $\mathrm{maxRank}(X') < \mathrm{maxRank}(X)$.
\end{lemma}

\section{Decompositions}

\begin{definition}[Below and above]
  \label{def:below_above}
  \lean{Chromosome.below, Chromosome.above}
  \leanok
  \uses{def:Chromosome}
  For $k \in \mathbb{N}$, $X_{\le k}$ (below) filters genes with rank $\le k$,
  and $X_{> k}$ (above) filters genes with rank $> k$.
  We have $X = X_{\le k} + X_{> k}$.
\end{definition}

\begin{lemma}[Rank decomposition and lifting]
  \label{lem:Chromosome.rank_decomposition}
  \lean{Chromosome.rank_decomposition, Chromosome.prime_iterate_eq_prime_iterate_above, Chromosome.lift_prime}
  \leanok
  \uses{def:below_above, def:Chromosome.prime, def:Chromosome.lift}
  The below/above filters decompose a chromosome.  The below part of rank
  at most $k$ is killed by $k$ primes, so $X$ can be reconstructed from
  $X^{(k)}$ by lifting, together with its below-$k$ part.
\end{lemma}

\begin{definition}[Parity decomposition]
  \label{def:parity}
  \lean{Chromosome.oddPart, Chromosome.evenPart}
  \leanok
  \uses{def:Chromosome}
  Every chromosome decomposes as $X = X_{\mathrm{odd}} + X_{\mathrm{even}}$
  where $X_{\mathrm{odd}}$ (resp.\ $X_{\mathrm{even}}$) collects genes of odd (resp.\ even) rank.
\end{definition}

\begin{lemma}[Parity and prime interchange]
  \label{lem:parity_prime}
  \lean{Chromosome.evenPart_prime, Chromosome.oddPart_prime}
  \leanok
  \uses{def:parity, def:Chromosome.prime}
  $(X')_{\mathrm{even}} = (X_{\mathrm{odd}})'$ and $(X')_{\mathrm{odd}} = (X_{\mathrm{even}})'$.
\end{lemma}

\begin{remark}[Module organization]
  The Lean development now follows the mathematical blocks Basic, Signature,
  Prime, Rank, Order, Lift, and Parity under the \texttt{Chromosome}
  directory.  The module \texttt{YoungDiagram.Chromosome} re-exports them as
  the stable public entry point.
\end{remark}

% ================================================================

\chapter{Varieties}

\section{Filtered varieties}

\begin{definition}[Variety]
  \label{def:Variety}
  \lean{Variety}
  \leanok
  \uses{def:Chromosome}
  A \emph{variety} is an additive submonoid of $\mathtt{Chromosome}$.
\end{definition}

\begin{definition}[Filtered variety]
  \label{def:varietyOfFilter}
  \lean{Chromosome.varietyOfFilter}
  \leanok
  \uses{def:Variety}
  Given a predicate $p$ on genes, the \emph{filtered variety} $V_p$ consists of
  all chromosomes whose support satisfies $p$.
\end{definition}

\begin{lemma}[Filtered varieties are prime-stable]
  \label{lem:prime_varietyOfFilter}
  \lean{Chromosome.prime_varietyOfFilter}
  \leanok
  \uses{def:varietyOfFilter, def:Chromosome.prime}
  If $p$ is lift-stable (i.e.\ $p(g)$ iff $p(\mathrm{lift}(g))$), then the prime
  of a filtered variety is itself: $V_p' = V_p$.
\end{lemma}

\section{The polarized variety Pi}

\begin{definition}[Polarized chromosome]
  \label{def:IsPolarized}
  \lean{Chromosome.IsPolarized}
  \leanok
  \uses{def:Chromosome}
  A chromosome $X$ is \emph{polarized} if every gene in its support has type
  $\ne \mathtt{NonPolarized}$.
\end{definition}

\begin{definition}[Variety Pi]
  \label{def:Variety.Pi}
  \lean{Variety.Pi}
  \leanok
  \uses{def:varietyOfFilter, def:IsPolarized}
  The variety $\mathrm{Pi}$ is the filtered variety of all polarized chromosomes.
\end{definition}

\begin{lemma}[Pi is prime-stable]
  \label{lem:prime_Pi}
  \lean{Variety.prime_Pi}
  \leanok
  \uses{def:Variety.Pi}
  $\mathrm{Pi}' = \mathrm{Pi}$.
\end{lemma}

\begin{lemma}[Polarized signature is natural]
  \label{lem:signature_pi_isNat}
  \lean{signature_pi_isNat}
  \leanok
  \uses{def:Variety.Pi, def:Chromosome.signature}
  If $X \in \mathrm{Pi}$, the signature components of $X$ are natural numbers.
\end{lemma}

\begin{lemma}[Sigma sequence determines element of Pi]
  \label{lem:eq_of_sigma_eq}
  \lean{sigmaUnique_Pi}
  \leanok
  \uses{def:Variety.Pi, def:Chromosome.signature, def:Chromosome.prime}
  If $X, Y \in \mathrm{Pi}$ satisfy $\mathrm{sig}(X^{(k)}) = \mathrm{sig}(Y^{(k)})$ for all $k \ge 0$,
  then $X = Y$.
\end{lemma}

\begin{lemma}[Antisymmetry of dominance on Pi]
  \label{lem:pi_chromosome_antisymm}
  \lean{Variety.SigmaUnique.partialOrder, sigmaUnique_Pi}
  \leanok
  \uses{def:Variety.Pi, def:Chromosome.Dominates}
  For $X, Y \in \mathrm{Pi}$, if $X \le Y$ and $Y \le X$, then $X = Y$.
  Hence $\mathrm{Pi}$ carries a partial order.
\end{lemma}

\section{The non-polarized variety Lambda}

\begin{definition}[Non-polarized chromosome]
  \label{def:IsNonPolarized}
  \lean{Chromosome.IsNonPolarized}
  \leanok
  \uses{def:Chromosome}
  A chromosome is \emph{non-polarized} if every gene in its support has type
  $\mathtt{NonPolarized}$.
\end{definition}

\begin{definition}[Variety Lambda]
  \label{def:Variety.Lambda}
  \lean{Variety.Lambda}
  \leanok
  \uses{def:varietyOfFilter, def:IsNonPolarized}
  The variety $\mathrm{Lambda}$ is the filtered variety of all non-polarized chromosomes.
\end{definition}

\begin{lemma}[Lambda is prime-stable]
  \label{lem:prime_Lambda}
  \lean{Variety.prime_Lambda}
  \leanok
  \uses{def:Variety.Lambda}
  $\mathrm{Lambda}' = \mathrm{Lambda}$.
\end{lemma}

\section{Mixed varieties and labels}

\begin{definition}[Mixed variety]
  \label{def:Variety.Mix}
  \lean{Variety.Mix}
  \leanok
  \uses{def:Variety}
  Given a pair of varieties $(V_1, V_2)$, the \emph{mixed variety} $\mathrm{Mix}(V_1, V_2)$
  consists of all chromosomes whose even part lies in $V_1$ and odd part lies in $V_2$.
\end{definition}

\begin{lemma}[Prime of a mixed variety]
  \label{lem:Variety.prime_Mix_eq}
  \lean{Variety.prime_Mix_eq}
  \leanok
  \uses{def:Variety.Mix, lem:parity_prime}
  If the two component varieties are closed under taking odd and even parts,
  then prime swaps the two entries of a mixed variety:
  $\mathrm{Mix}(V_1,V_2)'=\mathrm{Mix}(V_2',V_1')$.
\end{lemma}

\begin{definition}[Five labeled varieties]
  \label{def:Variety.Label}
  \lean{Variety.Label}
  \leanok
  \uses{def:Variety.Pi, def:Variety.Lambda, def:Variety.Mix}
  The five labeled varieties $V_0, \ldots, V_4$ indexed by $\mathrm{Fin}\,5$ are:
  \begin{itemize}
    \item $V_0 = \mathrm{Pi}$
    \item $V_1 = \mathrm{Mix}(\mathrm{Lambda}, \mathrm{Pi})$
    \item $V_2 = \mathrm{Mix}(\mathrm{Pi}, \mathrm{Lambda})$
    \item $V_3 = \mathrm{Mix}(2\mathrm{Lambda}, \mathrm{Pi})$
    \item $V_4 = \mathrm{Mix}(\mathrm{Pi}, 2\mathrm{Lambda})$
  \end{itemize}
\end{definition}

\begin{definition}[Label prime permutation]
  \label{def:Label.prime}
  \lean{Variety.Label.prime}
  \leanok
  \uses{def:Variety.Label}
  The prime operation on varieties induces a permutation of the five labels:
  $V_i' = V_{\pi(i)}$.
\end{definition}

\begin{lemma}[Prime tracks the labels]
  \label{lem:Variety.Label.prime_eq}
  \lean{Variety.Label.prime_eq, Variety.Label.prime_eq_iterate, Variety.prime_iterate_mem}
  \leanok
  \uses{def:Label.prime, lem:Variety.prime_Mix_eq}
  For every label $i$, the variety prime of $V_i$ is exactly
  $V_{\pi(i)}$.  Iterating this equality tracks membership after iterated
  prime operations.
\end{lemma}

% ================================================================

\chapter{Sigma Sequences}

\begin{definition}[Sigma sequence]
  \label{def:Sigma.sigma}
  \lean{Sigma.sigma}
  \leanok
  \uses{def:Chromosome.signature, def:Chromosome.prime}
  For a chromosome $X$, the \emph{sigma sequence} is the function
  $\sigma_X(k) = \mathrm{sig}(X^{(k)}) \in \mathbb{Q} \times \mathbb{Q}$.
  We write $a_X(k) = \sigma_X(k).1$ and $b_X(k) = \sigma_X(k).2$.
\end{definition}

\begin{lemma}[Sigma is antitone]
  \label{lem:Sigma.antitone}
  \lean{Sigma.antitone}
  \leanok
  \uses{def:Sigma.sigma}
  The sigma sequence is antitone (componentwise decreasing).
\end{lemma}

\begin{lemma}[First two sigma monotonicity conditions]
  \label{lem:Sigma.cond_15_2_3}
  \lean{Sigma.cond_15_2, Sigma.cond_15_3}
  \leanok
  \uses{def:Sigma.sigma, lem:Sigma.antitone, lem:Sigma.eventually_zero}
  The first and second components separately form decreasing sequences and
  both are eventually zero.  These are the formal versions of conditions
  (15.2) and (15.3).
\end{lemma}

\begin{lemma}[Sigma is eventually zero]
  \label{lem:Sigma.eventually_zero}
  \lean{Sigma.eventually_zero}
  \leanok
  \uses{def:Sigma.sigma}
  There exists $K$ such that $\sigma_X(k) = 0$ for all $k \ge K$.
\end{lemma}

\begin{lemma}[Interlacing conditions]
  \label{lem:Sigma.interlacing}
  \lean{Sigma.cond_15_4, Sigma.cond_15_5}
  \leanok
  \uses{def:Sigma.sigma}
  The components of the sigma sequence satisfy interlacing inequalities.
  If $k$ is even, then $b_X(k+1) \le a_X(k)$ and $a_X(k+1) \le b_X(k)$.
  If $k$ is odd, the inequalities are reversed.
\end{lemma}

\begin{lemma}[Difference conditions for Pi]
  \label{lem:Sigma.cond_15_6}
  \lean{Sigma.cond_15_6, Sigma.cond_15_7}
  \leanok
  \uses{def:Sigma.sigma, def:Variety.Pi}
  If $X \in \mathrm{Pi}$, the successive drops
  $a_X(k)-a_X(k+1)$ and $b_X(k)-b_X(k+1)$ satisfy the alternating
  inequalities (15.6) and (15.7).
\end{lemma}

\begin{lemma}[Sigma drops count alternating signs]
  \label{lem:Sigma.sigma_diff}
  \lean{Sigma.sigma_fst_diff, Sigma.sigma_snd_diff}
  \leanok
  \uses{def:Sigma.sigma, def:Variety.Pi}
  For a polarized chromosome, the drop in the first sigma component from
  $k$ to $k+1$ counts genes in $X^{(k)}$ that are positive in the alternating
  basis; the drop in the second component counts the corresponding negative
  genes.
\end{lemma}

\begin{lemma}[Comparing drops to the initial drop]
  \label{lem:Sigma.compare_drops}
  \lean{Sigma.cond_15_6_compare_drop_to_0, Sigma.cond_15_6_compare_k_to_0, Sigma.a1_ai_le_b0_bi_1}
  \leanok
  \uses{lem:Sigma.cond_15_6}
  The alternating inequalities imply useful global comparisons: later drops
  are bounded by the initial drop, and telescoping gives
  $a_1-a_i \le b_0-b_{i-1}$ for $i\ge 1$.
\end{lemma}

\begin{lemma}[Rank-one sign tests]
  \label{lem:Sigma.case_sign_tests}
  \lean{Sigma.neg_type_of_b0_gt_a1_single, Sigma.pos_type_of_b0_le_a1_single, Sigma.neg_gene_of_b0_gt_a1}
  \leanok
  \uses{def:Sigma.sigma, def:Variety.Pi}
  The inequalities between $b_0$ and $a_1$ detect the sign of a rank-one
  polarized gene.  Summing this test over a chromosome produces a negative
  gene whenever $b_0>a_1$.
\end{lemma}

\begin{lemma}[Dominance implies componentwise inequality]
  \label{lem:Sigma.cond_15_8}
  \lean{Sigma.cond_15_8}
  \leanok
  \uses{def:Sigma.sigma, lem:pi_chromosome_antisymm}
  If $X < Y$ in $\mathrm{Pi}$, then $\sigma_X(k) \le \sigma_Y(k)$ componentwise for all $k$
  with strict inequality for some $k$.
\end{lemma}
